\documentclass[12pt]{article}
\usepackage[utf8]{inputenc}
\usepackage[T1]{fontenc}
\usepackage{hyperref}
\title{Tasmanian, Southeast Language Resources for K-12 Curriculum Development}
\author{Compiled by Claude Sonnet 5.5 (claude-sonnet-5-5) for Andr\'as Kornai}
\date{}
\begin{document}
\maketitle

\section{Introduction}
Southeast Tasmanian (ISO 639-3 xpf, Glottocode sout1439) is an extinct language with very thin records. The searches found no source dedicated to it, no primer and no curriculum. The only living Tasmanian language is palawa kani, a reconstructed composite built by the Tasmanian Aboriginal Centre. It draws mainly on the Northeastern and Eastern Tasmanian families, not specifically on this language. Any heritage curriculum must be led by the Tasmanian Aboriginal community.

\section{Primary Educational Resources}
None specific to this language were found. The palawa kani programme is described at \url{https://tacinc.com.au/programs/palawa-kani/}. The Centre reportedly allows outside access only to place names and gives dictionaries and learning materials only to the Aboriginal community.

\section{Linguistic and Reference Materials}
\begin{itemize}
\item Overview of the Tasmanian languages and their word lists: \url{https://en.wikipedia.org/wiki/Tasmanian_languages}.
\item N.\,J.\,B. Plomley's 1976 word-list compilation, a primary source, is catalogued at \url{https://glottolog.org/resource/reference/id/20405}.
\item The Oxford Guide to Australian Languages chapter on Tasmanian is at \url{https://academic.oup.com/book/46706/chapter/410243287}.
\end{itemize}

\section{Religious and Community Resources}
The Tasmanian Aboriginal Centre is the community body responsible for palawa kani: \url{https://tacinc.com.au/programs/palawa-kani/}. Search results describe the only continuous text as a 1829 sermon recorded by G.~A. Robinson on Bruny Island, and note that it is largely English with words swapped. No other religious material was found.

\section{Digital and Audio Resources}
The search results mention a 1899 recording of Fanny Cochrane Smith singing traditional songs, discussed at \url{https://en.wikipedia.org/wiki/Palawa_kani}. No archive link was found in search results.

\section{Curriculum Development Context}
The evidence for Southeast Tasmanian is a few word lists only. Immersion or full-language teaching is not possible. A heritage unit would most plausibly be about the history and languages of Tasmania, with palawa kani taught under the community's own terms.

\end{document}